\chapter*{Preface}
\addcontentsline{toc}{chapter}{Preface}
\label{chap:preface}

This dissertation was not written from inside an institution. There is no advisor whose approval shaped the research direction, no committee whose preferences constrained the questions worth asking, no department whose orthodoxies set the boundaries of what could be claimed. The work originated in 2016 as a self-funded experiment in language-model behavior, conducted without grant support, without laboratory infrastructure, and without any expectation that it would eventually produce a research program spanning thirty-four projects, a formal algebraic equation, and a doctoral-level argument about the failure mode of contemporary artificial intelligence deployment.

That independence is not a credential gap. It is the most important structural fact about what follows.

Independent research is not research conducted in ignorance of the literature. Every claim in this corpus cites a source: a paper, an experiment, a checkpoint, or a prior finding. The Van der Aalst OCEL 2.0 specification grounds the process-evidence layer. The Object-Centric Event Log standard grounds the replay claims. The cryptographic receipt chain in \texttt{receipts/} grounds the provenance claims. What independent research does is eliminate the social incentive to fabricate closure. There is no funding cycle that rewards a premature \textsc{alive} verdict. There is no publication queue that benefits from overstating precision. The research program documented here operates under a strict refusal law: if the receipt chain does not confirm a claim, the claim is \textsc{partial}, not complete.

\medskip

This dissertation is customer-first, not investor-first.

That distinction is not rhetorical. It is structural. Every chapter of this work is organized around a single question: can a customer --- a hospital, a law firm, a capital-markets desk, a municipality --- verify that consequential work occurred? Not: can an investor be shown a benchmark that implies it might occur. Not: can a demo be constructed that simulates the appearance of it occurring. The question is whether the work actually happened, whether that happening is provable, and whether the proof is durable.

The difference between those two framings determines everything that follows. Investor-first logic selects for language that sounds like progress. Customer-first logic selects for receipts that prove progress. This dissertation is organized around receipts.

\medskip

The title of the dissertation is \textit{From Prediction to Receipted Coordination}. That phrase describes the exact arc of the argument. Language models predict. That is a useful capability. But prediction is not coordination. A prediction does not hand off to the next process stage. A prediction does not produce a receipt that can be replayed. A prediction does not close work. The movement from prediction to receipted coordination is the movement from a capability that impresses observers to a capability that serves customers.

This work documents that movement. It documents it through a 2016 experiment that revealed the failure mode. It documents it through the Chatman Equation, which provides the algebraic formalization. It documents it through ggen, which manufactures artifacts under law. It documents it through the Open Ontologies program, which makes the meaning-bearing layer public and auditable. It documents it through Post-Cyberpunk PCP, which names the cultural moment in which hallucination-as-output is being displaced by receipt-as-proof. It documents it through AI XYNZ and the capital-flow infrastructure, which connects the equation to the securities authority stack. And it documents it through thirty-four discovered projects, each with its own \textsc{alive}/\textsc{partial}/\textsc{blocked} verdict, forming a corpus of evidence that can be mined, replayed, and audited.

\medskip

Why is this a crescendo?

Because the 2016 experiment was a beginning, not a conclusion. The observation that local text imitation does not conserve consequence was a single data point. The decade of work since --- the process-gap analysis, the equation, the manufacturing machinery, the open ontologies, the receipt chain --- has converted that observation into an operating architecture. The crescendo is not the architecture itself. The crescendo is the proof that the architecture works: thirty-four projects, a verifiable receipt chain, conformance-checked process logs, and a customer-first evaluation that does not rely on investor-facing metrics.

The argument is complete. The receipts exist. The work is done.

\vspace{2em}
\noindent\textit{Sean Chatman}\\
\noindent\textit{Process Intelligence Research Foundry}\\
\noindent\textit{Independent, 2026}
